
\subsection{$\mathcal M_{k+1}^{\mathrm{main}}(\beta_0)$
的仓西结构}
\label{subsec:7-5-3}

为应用第~\ref{subsec:7-5-2} 小节所发展的方法，
我们回顾
$\mathcal M_{k+1}^{\mathrm{main}}(\beta_0)$ 与
$\mathcal M_{k+1}^{\mathrm{main},\mathfrak d}
(\beta_0;\vec P)$
的仓西结构的显式描述。
回忆 $\mathcal M_{k+1}^{\mathrm{main}}$ 表示
边界上带 $k+1$ 个标记点的圆盘复结构模空间的主分支
（见定义~\ref{def:2-1-18}）。
于是，作为空间，
$\mathcal M_{k+1}^{\mathrm{main}}(\beta_0)$
同 $L\times\mathcal M_{k+1}^{\mathrm{main}}$ 识别，
而评价映射同对角映射
\[
L\longrightarrow L^{k+1},
\qquad
p\longmapsto(p,\ldots,p)
\]
识别。
由于评价映射不是浸没，
需要引入障碍丛并扩张仓西邻域
（见第~\ref{sec:3-4} 节）。
具体地，设 $U_{k+1}$ 是 $L$ 在 $L^{k+1}$ 中的一个邻域；
这里把 $L$ 看成 $L^{k+1}$ 的大对角线。
把 $U_{k+1}$ 同 $L\subset L^{k+1}$ 的法丛
$N_L L^{k+1}$ 识别，并置
\[
U_{k+1}(\mathfrak m)
=U_{k+1}\times\mathcal M_{k+1}^{\mathrm{main}}.
\]

设 $E_{k+1}$ 是法丛 $N_L L^{k+1}$
经投影 $U_{k+1}\to L$ 拉回到 $U_{k+1}$ 所得的丛，
并把它在 $U_{k+1}(\mathfrak m)$ 上的拉回记作
$E_{k+1}(\mathfrak m)$。
识别 $U_{k+1}\subset N_L L^{k+1}$ 诱导截面
\[
s_{\mathfrak m,k+1}\colon
U_{k+1}(\mathfrak m)\longrightarrow E_{k+1}(\mathfrak m),
\qquad
s_{\mathfrak m,k+1}^{-1}(0)
=L\times\mathcal M_{k+1}^{\mathrm{main}}.
\]
注意，$s_{\mathfrak m,k+1}$ 关于第二个因子
$\mathcal M_{k+1}^{\mathrm{main}}$ 是常值的。
取
\[
\bigl(U_{k+1}(\mathfrak m),
E_{k+1}(\mathfrak m),s_{\mathfrak m,k+1}\bigr)
\]
作为 $\mathcal M_{k+1}^{\mathrm{main}}(\beta_0)$
的仓西结构
\translatornote{源书此处漏去上标 $\mathrm{main}$，写作
$\mathcal M_{k+1}(\beta_0)$；但小节标题、开篇所列对象及
紧前零集 $L\times\mathcal M_{k+1}^{\mathrm{main}}$
都明确限定在主分支，故译文补回。见疑误 A-0616。}。
由于 $U_{k+1}$ 是 $L^{k+1}$ 的开子集，
评价映射
\[
ev\colon U_{k+1}(\mathfrak m)\longrightarrow L^{k+1}
\]
是浸没。
还要指出，自同构群是平凡的。

从这里起，我们都使用
$U_{k+1}(\mathfrak m)$、$E_{k+1}(\mathfrak m)$、
$s_{\mathfrak m,k+1}$ 等记号。
它们定义用来构造算子 $\mathfrak m$ 的仓西结构
\[
\bigl(U_{k+1}(\mathfrak m),
E_{k+1}(\mathfrak m),s_{\mathfrak m,k+1}\bigr).
\]
对于用来构造 $\mathfrak f$ 的模空间，
下一小节将写作
\[
\bigl(U_{k+1}(\mathfrak f),
E_{k+1}(\mathfrak f),s_{\mathfrak f,k+1}\bigr).
\]
存在评价映射
\[
ev=(ev_1,\ldots,ev_k,ev_0)\colon
U_{k+1}(\mathfrak m)\longrightarrow L^{k+1}.
\]
我们用字母 $s$（如 $s_{\mathfrak m,k+1}$）
表示原仓西映射，用 $\mathfrak s$ 表示它的扰动（截面），
用 $\mathfrak t$ 表示它的连续扰动族。

设 $P_i\in\mathcal X_L$，并以 $\mathfrak d(i)$ 表示它的代。
置 $\vec P=(P_1,\ldots,P_k)$。
下面描述
$\mathcal M_{k+1}^{\mathrm{main},\mathfrak d}
(\beta_0;\vec P)$ 的仓西结构。
按定义，其仓西邻域是纤维积
\[
U_{k+1}(\mathfrak m)
\mathbin{\times_{L^k}}(P_1\times\cdots\times P_k).
\]
它按定义同下式右端同构：
\begin{equation}
U_{k+1}(\mathfrak m,\vec P)
=\bigl(U_{k+1}\times\mathcal M_{k+1}^{\mathrm{main}}\bigr)
\mathbin{\times_{L^{k+1}}}
(P_1\times\cdots\times P_k\times L).
\tag{7.5.12}\label{eq:7-5-12}
\end{equation}
这里要指出，在 $L^{k+1}$ 中，
最后一个分量是第 $0$ 个标记点处评价映射的靶。
截面 $s_{\mathfrak m,k+1}^{\vec P}$ 是
$E_{k+1}(\mathfrak m)$ 在
$U_{k+1}(\mathfrak m,\vec P)$ 上的限制丛的截面。
